\section{Discusión}
\label{sec:discusion}

Los resultados obtenidos ponen de manifiesto por qué un dimensionamiento basado únicamente en promedios diarios o en fórmulas de régimen estacionario resulta inadecuado cuando la tasa de arribos varía a lo largo del día.
En los tres escenarios analizados, las dotaciones recomendadas ($c^* = 3$, 4 y 6) operan con utilizaciones medias sobre las 13~horas de apertura moderadas, del 65,5~\% en demanda baja, 58,4~\% en demanda media y 63,2~\% en demanda alta (Tabla~\ref{tab:experimentos-escenarios}).
Una lectura estática sugeriría que existe capacidad ociosa suficiente para reducir un preparador; sin embargo, al pasar a $c^* - 1$ empleados ($c = 2$, 3 y 5), la utilización media asciende al rango de 73,7~\% a 89,7~\% y el cumplimiento del plazo de 30~minutos cae abruptamente a 37,5~\%, 74,0~\% y 84,0~\%.
Esto ocurre porque el perfil horario concentra la demanda en dos picos pronunciados alrededor de las 11 y las 19~horas: durante esas franjas la tasa instantánea de llegadas alcanza 18,9, 21,9 y 38,2~pedidos por hora en demanda baja, media y alta (Tabla~\ref{tab:experimentos-recomendacion}), superando la capacidad media de atención de $c = 2$, 3 y 5~preparadores, que con 10~minutos de armado por pedido en promedio es de 12, 18 y 30~pedidos por hora.
Incluso con la dotación recomendada, en demanda baja y alta el pico supera la capacidad media de $c^* = 3$ y 6~preparadores (18 y 36~pedidos por hora), por lo que la fila crece durante la hora pico y se vacía en las horas siguientes, de menor demanda.
La diferencia está en el arrastre: con $c^* - 1$~preparadores el exceso acumulado en el pico tarda más en drenarse, alcanza a los pedidos de las horas siguientes y deteriora la proporción global de entregas a tiempo, aun cuando en las primeras horas de la mañana o en la siesta el personal permanezca parcialmente desocupado.

En segundo lugar, la articulación entre la etapa de pronóstico y la simulación permite cuantificar el impacto operativo de los errores de predicción.
La evaluación en el período de prueba mostró que la media estacional de ocho semanas presenta un error absoluto medio de 3,42~pedidos por hora (muy próximo al piso de 3,21~pedidos por hora del proceso de Poisson subyacente) y un sesgo de subestimación del 4,4~\% producto de la tendencia creciente del histórico (Tabla~\ref{tab:errores-pronostico}).
Al propagar este pronóstico hacia el simulador y contrastar la decisión sobre las cuatro semanas reservadas de prueba (Tabla~\ref{tab:experimentos-recomendacion}), se verifica que la naturaleza discreta de la dotación introduce un margen de absorción desigual entre escenarios.
En demanda media ($c^* = 4$) y alta ($c^* = 6$), el salto discreto respecto de $c^* - 1$ deja un cumplimiento nominal de 97,4~\% y 97,5~\%, lo que permite absorber sin cambios de personal tanto el sesgo observado en los días de prueba (100,0~\% y 98,1~\% de cumplimiento) como una subestimación sistemática del 10~\% en la sensibilidad (94,3~\% y 93,2~\% en la Tabla~\ref{tab:sensibilidad-recomendacion}).
En cambio, en el escenario de demanda baja la dotación $c^* = 3$ queda cerca del límite exigido (92,0~\% en la simulación nominal y 90,6~\% en los domingos de prueba), por lo que un incremento del 10~\% en las llegadas basta para reducir el cumplimiento al 85,3~\% y requerir un cuarto preparador.

En tercer término, el estudio de sensibilidad evidencia que la incertidumbre sobre los parámetros del tiempo de servicio ($b$ y $m$) tiene un peso tanto o más determinante que el error del pronóstico horario.
Mientras que la diferencia entre el pronóstico medio y un aumento del 10~\% en la demanda altera la dotación óptima únicamente en el escenario bajo, un incremento del 20~\% en el tiempo medio de armado (pasar de 10 a 12~minutos por pedido) obliga a sumar un preparador en todos los escenarios ($c^* = 4$, 5 y 7), y un incremento del 40~\% (14~minutos por pedido) requiere dos preparadores adicionales ($c^* = 5$, 6 y 8).
Dado que el tiempo de servicio multiplica a cada uno de los pedidos que ingresan al sistema y reduce simultáneamente la tasa de drenaje de las filas acumuladas en las horas pico, cualquier fricción operativa no registrada en el ERP ---como demoras en la búsqueda de artículos, congestión en pasillos o verificación de productos frescos--- impacta de manera directa sobre la capacidad efectiva del sector.

Finalmente, la comparación de los tres escenarios demuestra que una dotación única e invariable para toda la semana resulta ineficiente.
Fijar una plantilla permanente de $c = 6$~preparadores todos los días permitiría cubrir los sábados de demanda alta con un 97,5~\% de cumplimiento, pero en días de demanda baja y media dejaría al personal con una utilización del 33,4~\% y 39,4~\%, respectivamente (Tabla~\ref{tab:experimentos-escenarios}).
Por el contrario, fijar una dotación uniforme de $c = 4$~preparadores cubriría adecuadamente los días bajos y medios, pero provocaría el colapso del servicio los sábados, con apenas un 50,8~\% de pedidos a tiempo, una espera media de 26,74~minutos y 31,3~pedidos pendientes al cierre del local.
Adaptar la dotación según el nivel de demanda pronosticado para cada tipo de jornada constituye, por lo tanto, una regla operativa simple y de alto impacto para equilibrar calidad de servicio y uso de recursos.
